\documentclass[11pt,letterpaper]{article}
\usepackage[T1]{fontenc}
\usepackage{lmodern,amsmath,amssymb,amsthm,microtype}
\usepackage[margin=1in]{geometry}
\usepackage[hidelinks]{hyperref}
\newtheorem{lemma}{Lemma}
\newtheorem{theorem}{Theorem}
\title{Five isolates are necessary in a minimum-ground counterexample to Frankl's conjecture at height six}
\author{Aafi Khan Malek\\ \small Independent researcher, Ahmedabad, India\\
\small\href{mailto:aafimalek2032@gmail.com}{aafimalek2032@gmail.com}}
\date{28 November 2026}
\begin{document}
\maketitle
Use the weak convention: a counterexample is a nonempty finite
union-closed family of nonempty sets with no element in more than
half its members. Height is the maximum number of members in a
strict inclusion chain. Ground minimality here means minimum ground
size among weak counterexamples of height at most six. A point
forces a distinct point if every set containing the first contains
the second; a point is isolated if there is no such comparison
to or from it.

\begin{theorem}
A weak ground-minimal counterexample of height at most six has at
least five isolated points.
\end{theorem}
This is a structural reduction, not a proof of the unrestricted
height-six problem. The repair argument below is written mathematics
and the inherited finite trace step is reproducible, but the theorem
has not yet been checked end to end in Lean.

\section{Repair by maximal missing deletion}
For a family \(\mathcal B\), let \(f_{\mathcal B}(v)\) count the
members containing \(v\).
\begin{lemma}
Let \(\mathcal B\) be a nonempty union-closed subfamily of a
finite union-closed family \(\mathcal G\) of nonempty sets. Suppose
\(\mathcal G\) has height at most six and some member strictly
contains every member of \(\mathcal B\). Put \(b=|\mathcal B|\).
Suppose \(\{x\}\notin\mathcal B\),
\[
 d=2f_{\mathcal B}(x)-b>0,\qquad
 2f_{\mathcal B}(v)\le b-d\quad(v\ne x).
\]
Then there is a weak counterexample of height at most six on the
ground of \(\mathcal B\).
\end{lemma}

\begin{proof}
Write \(\mathcal B_0=\mathcal B\). For \(0\le t<d\), the difference
between the numbers of \(x\)-containing and \(x\)-avoiding members of
\(\mathcal B_t\) is \(d-t>0\), since each of the previous \(t\) insertions
avoids \(x\). Deletion of \(x\) is injective on the members containing
it. Hence there is such a member \(A\) with missing deletion
\(C=A\setminus\{x\}\). Choose \(A\) inclusion-maximal among these members
and put \(\mathcal B_{t+1}=\mathcal B_t\cup\{C\}\). Every chosen \(A\)
is in the original \(\mathcal B\), and \(C\ne\varnothing\).

For \(D\in\mathcal B_t\) containing \(x\), we have \(C\cup D=A\cup D\)
in \(\mathcal B_t\). Otherwise \(C\cup D=(A\cup D)\setminus\{x\}\).
If \(A\cup D=A\) this is the new member \(C\); if \(A\cup D\) is
strictly larger, it belongs to \(\mathcal B_t\) and its deletion
already belongs there by maximality. This proves union closure
inductively. At \(t=d\) the size is \(b+d\), the unchanged
\(x\)-frequency is \((b+d)/2\), and the doubled frequency of any
other point is at most \(b-d+2d=b+d\). The ground is unchanged.

It remains to prove the height bound directly for \(\mathcal B_d\).
Take a strict chain using at least one inserted member; let its first
inserted member be \(C=A\setminus\{x\}\) with \(A\in\mathcal B\).
Leave the prefix below \(C\) unchanged. Add \(x\) to each set in
the suffix beginning at \(C\). An original suffix member \(D\)
contains \(C\), so its lift is \(A\cup D\in\mathcal B\);
a newly inserted suffix member lifts to its own original preimage
in \(\mathcal B\). The boundary below \(C\) remains strict because
the prefix omits \(x\) and the lifted \(C\) contains it.
Adjoining \(x\) can identify adjacent distinct suffix members
only at the step that adds \(x\) and nothing else. There is at most
one such step along a chain. If it is lost, append the member of
\(\mathcal G\) strictly above \(\mathcal B\); this restores the
length. Thus every such chain lifts to one of the same length
in \(\mathcal G\). Chains with no inserted members are already in
\(\mathcal G\), proving the desired height bound.
\end{proof}

\section{Application to four isolates}
The \href{https://github.com/Aafimalek/novel-discoveries/tree/2d50235bd73e18b1c1371f96dd7fc3540514e0ee/frankl-height-six/verification/at-least-four-isolates}{at-least-four reduction}
excludes fewer than four isolates. At exactly four, the previous
\href{https://github.com/Aafimalek/novel-discoveries/tree/931f364753a6e70cc6b093e09b42e473134beb04/frankl-height-six/verification/exactly-four-isolates}{block and trace reductions}
leave only a single forcing source \(p\). The containing fiber
\(\mathcal P\) consists of the full forcing star together with
all 15 isolate traces except one singleton \(\{x\}\).
The finite step examines 106 eligible labeled families and
excludes 105; \texttt{python3 traces.py} reproduces it.

Let \(\mathcal B\) consist of members not containing \(p\).
It is union-closed, on smaller ground, and nonempty since
\(p\) cannot be universal. The union maximum of \(\mathcal G\)
contains \(p\) and is strictly above \(\mathcal B\).
In \(\mathcal P\), \(x\) occurs seven times, each other isolate
eight times, and each leaf fifteen times. The weak frequency bounds
in \(\mathcal G\) therefore imply
\[
 2f_{\mathcal B}(x)\le b+1,\quad
 2f_{\mathcal B}(y)\le b-1\ \text{for other isolates},\quad
 2f_{\mathcal B}(\ell)\le b-15\ \text{for leaves}.
\]
By ground minimality \(\mathcal B\) must have a strict majority.
Only \(x\) can provide it, forcing \(d=1\). Also \(\mathcal B\)
does not contain the singleton \(\{x\}\). In a weak counterexample
of nonempty sets, insertion of a member singleton into all
avoiding members would inject them into the containing members
other than the singleton itself, contradicting the frequency bound.
The lemma produces a weak counterexample on smaller ground, a
contradiction. This excludes exactly four isolates and proves
the theorem.

\paragraph{Further scope.}
For any one-source containing fiber \(\mathcal P\), the same repair
applies if there is an isolate \(x\) with
\(2f_{\mathcal P}(x)<|\mathcal P|\) and every other active point
has \(f_{\mathcal P}(v)\ge |\mathcal P|-f_{\mathcal P}(x)\).
Indeed the full-family weak bounds give at most
\(|\mathcal P|-2f_{\mathcal P}(x)\) surplus for \(x\) in
\(\mathcal B\) and at least that much deficit for each other
point. A strict majority in \(\mathcal B\) again has to be \(x\).
In particular, a full containing trace cube except for \(\{x\}\)
is excluded for every isolate count. This does not eliminate
all five-isolate profiles.

\paragraph{Related work.}
\href{https://arxiv.org/abs/2608.25147}{Tian} studies bounded height in
the empty-set-free convention. At height five,
\href{https://zenodo.org/records/23252932}{Wang's preprint} reports a
union-closure-preserving repair of a reduced family and
\href{https://zenodo.org/records/23244603}{Shimizu's preprint} reports a
Lean-checked theorem. These statements do not settle arbitrary height six
under the convention of this note. Our result is the application of a
quantitative deletion repair to the four-isolate reduction, not a
priority claim for the repair idea itself.

\medskip\noindent
This research note was developed with AI assistance.
\end{document}
